\section{Actor--Critic 在偏差与方差之间搭桥}
\label{sec:16-Control-and-Reinforcement-Learning/02-Reinforcement-Learning/10}

抓取任务，一个回合两百步，成败在最后一步才见分晓。REINFORCE 跑了三天，回报曲线上下翻飞，五个随机种子给出五种走势。你粗算了一下噪声：每步奖励的方差是 \(\sigma^2\)，折扣 \(\gamma=0.99\) 时回报的方差大致是 \(\sigma^2/(1-\gamma^2)\approx50\sigma^2\)，而每一步动作拿到的梯度权重就是这个抖动了五十倍的数。梯度无偏，可单个样本的符号都不太可信。

你想要的是在第 \(t\) 步当场给这个动作打分，不必等两百步之后。当场打分需要一个会预测未来的模型，模型一定有误差。于是选择摆在这里：等到底，估计正确但抖；当场打，估计稳定但偏。

\subsection{Critic 是策略梯度里被替换掉的那一项}

策略梯度定理可以写成

\[
\nabla_\theta J(\theta)=\E_\pi\big[Q^\pi(S_t,A_t)\,\nabla_\theta\log\pi_\theta(A_t\given S_t)\big],
\]

上一节用采样回报 \(G_t\) 当 \(Q^\pi\) 的无偏样本，方差有多大已经看到了。换个做法：学一个 \(V_\phi\approx V^\pi\)，用它就地构造权重。减去状态基线得到优势 \(A^\pi(s,a)=Q^\pi(s,a)-V^\pi(s)\)，而一步时序差分

\[
\delta_t=R_{t+1}+\gamma V_\phi(S_{t+1})-V_\phi(S_t)
\]

正是优势的一个估计。骨架一行：在 \(V_\phi=V^\pi\) 的前提下，对下一状态取期望有 \(\E[R_{t+1}+\gamma V^\pi(S_{t+1})\given S_t,A_t]=Q^\pi(S_t,A_t)\)，减去 \(V^\pi(S_t)\) 便得 \(A^\pi(S_t,A_t)\)。条件对账只有一条，但它很硬：\(V_\phi\) 得等于 \(V^\pi\)。Critic 有多少误差，Actor 的梯度就带多少偏差，而且这偏差不随样本量增加而消失。

于是两条更新交织进行，Actor 用

\[
\theta\leftarrow\theta+\alpha_\theta\,\delta_t\,\nabla_\theta\log\pi_\theta(A_t\given S_t),
\]

Critic 用半梯度 TD 最小化 \(\delta_t^2\)。同一条样本喂两个模型，一个学怎么做，一个学做得怎么样。

\subsection{一步与整条之间是一条连续的轴}

把真实奖励用满 \(n\) 步再接上估计，得到多步回报

\[
G_t^{(n)}=\sum_{k=0}^{n-1}\gamma^kR_{t+k+1}+\gamma^nV_\phi(S_{t+n}).
\]

\(n=1\) 是 TD(0)，\(n\) 取到回合结束就是 Monte Carlo。\(n\) 越大，进入估计量的真实奖励越多，对 \(V_\phi\) 的依赖越轻，偏差越小；同时被累加的随机步数越多，方差越大。Actor--Critic 的关键设计选择落在这个截断点的位置上，用了几个网络反倒是次要的。

\begin{center}
\begin{tikzpicture}[>=stealth,scale=0.9]
\fill[gray!22] (0,0.35) -- (8,0.35) -- (8,1.85) -- cycle;
\fill[gray!22] (0,2.3) -- (0,3.8) -- (8,2.3) -- cycle;
\draw[->] (-0.3,0) -- (8.8,0);
\foreach \x in {0,2.5,5,6.3,8} \draw (\x,-0.1) -- (\x,0.1);
\node[below,font=\small] at (0,-0.12) {\(\lambda=0\)};
\node[below,font=\small] at (2.5,-0.12) {\(0.5\)};
\node[below,font=\small] at (5,-0.12) {\(0.9\)};
\node[below,font=\small] at (6.3,-0.12) {\(0.95\)};
\node[below,font=\small] at (8,-0.12) {\(1\)};
\node[right,font=\small] at (5.6,1.0) {方差};
\node[right,font=\small] at (0.5,2.75) {偏差};
\node[above right,font=\small,align=left] at (-0.3,3.95) {TD(0)\\\(n=1\)};
\node[above left,font=\small,align=right] at (8.6,3.95) {Monte Carlo\\\(n=\infty\)};
\draw[dashed] (6.3,-0.1) -- (6.3,3.7);
\node[above,font=\small] at (6.3,3.7) {常用取点};
\end{tikzpicture}
\end{center}

\emph{一条轴，两个楔形：左端偏差最厚、方差最薄，右端相反；\(\lambda\) 与 \(n\) 只是这条轴上的两种刻度。}

图上 \(\lambda\) 的刻度是不均匀的，这不是画歪了。\(\lambda\) 起作用是通过有效视界 \(1/(1-\gamma\lambda)\)：\(\gamma=0.99\) 时，\(\lambda=0.9\) 对应大约十步，\(\lambda=0.95\) 对应二十步，\(\lambda=0.99\) 一下子跳到五十步。旋钮在右端极其灵敏，调参时把 \(\lambda\) 从 \(0.95\) 改成 \(0.99\) 与从 \(0.5\) 改成 \(0.9\) 完全不是同一量级的改动。

\subsection{GAE 把这条轴变成一个旋钮}

不必在 \(n=1\) 和 \(n=\infty\) 之间二选一。把各阶多步优势按 \((\gamma\lambda)^l\) 指数加权平均，展开后所有中间项抵消，只剩下 TD 残差的加权和：

\[
\widehat A_t^{\mathrm{GAE}(\lambda)}=\sum_{l=0}^{\infty}(\gamma\lambda)^l\,\delta_{t+l}.
\]

两端可以直接验算。\(\lambda=0\) 时只剩 \(\delta_t\)，回到一步 TD。\(\lambda=1\) 时

\[
\begin{aligned}
\sum_{l\ge0}\gamma^l\delta_{t+l}
&=\sum_{l\ge0}\gamma^l\big(R_{t+l+1}+\gamma V_\phi(S_{t+l+1})-V_\phi(S_{t+l})\big)\\
&=G_t-V_\phi(S_t),
\end{aligned}
\]

价值项逐级抵消，只留下首项，得到 Monte Carlo 优势。\(\lambda\) 在中间时，近处的残差权重大、远处按 \(\gamma\lambda\) 衰减，偏差与方差的交换点被一个连续参数定住。TD(\(\lambda\)) 的资格迹是同一件事在增量形式下的写法。

\subsection{两个学习器在互相追赶}

Critic 要评估的是``当前策略''的价值，而 Actor 每一步都在改动当前策略。Critic 的目标一直在走。

这带来一种独特的失效方式。Actor 跑得太快，Critic 手上的估计属于几十步之前那个策略，套到现在就是系统性偏差；Actor 依偏差方向走，策略再变，Critic 更追不上。两个误差可以互相喂养。理论分析通常要求两个时间尺度：Critic 的学习率显著大于 Actor，使它在 Actor 每走一步的时间里近似追上当前策略的价值。深度实现很难严格满足这类渐近条件，但这个视角解释了一批看似零散的技巧——更新频率比、相对学习率、Critic 侧的目标网络、延迟若干步再更新策略——它们调的是同一件事：两个学习器之间的相对节奏。

策略侧的信赖域和比率裁剪也可以放进这个视角看：限制单次策略改动的幅度，等于给 Critic 留出追赶的时间。

\subsection{离策略版本要额外付一份账}

数据来自行为策略 \(\mu\)、优化目标是 \(\pi_\theta\) 时，残差需要按

\[
\rho_t=\frac{\pi_\theta(A_t\given S_t)}{\mu(A_t\given S_t)}
\]

加权，多步展开则要乘上一串比率。上一节说过这串乘积会怎样集中，V-trace 之所以对两个阈值分开处理，正是因为其中一个决定不动点对应哪个策略，另一个控制迹沿轨迹传播的强度。在 Actor--Critic 里这道校正要做两次：Critic 的目标需要校正到 \(\pi_\theta\) 上，Actor 的梯度也需要，两处的截断偏差会叠加。

分布式训练里的滞后同样是这件事。采样器上跑的策略比学习器上的旧几十次更新，数据天然离策略，重要性权重的分布直接由这个滞后决定。把采样器与学习器的同步频率当作一个和 \(\lambda\) 同等重要的超参数，比事后调裁剪阈值有用。

\subsection{共享表示要划清梯度边界}

Actor 与 Critic 共用前几层能省参数和算力，两个目标的梯度却未必朝同一个方向：Critic 要求表示能准确预测价值，Actor 要求表示支持策略改进，两者有时互相促进，有时正面冲突。常见做法是给两侧的损失配上不同权重，或者干脆分成两个网络，用一点冗余换清晰。

比权重更要紧的是梯度边界。优势 \(\widehat A_t\) 是由 Critic 算出来的，它出现在 Actor 的损失里；忘了在这里切断梯度，Actor 的损失就会顺着优势回传进 Critic，把 Critic 往``让优势更大''的方向拉——一个谁都没打算优化的目标。哪些量是被学习的、哪些量是当作常数的目标，应当在代码里显式写出来，别交给框架的默认行为。

\subsection{多一个网络不改变底层的困难}

Actor--Critic 压下了策略梯度的方差，同时请进了 Critic 的逼近误差、自举带来的偏差，以及两个目标互相牵动的非平稳性；离策略版本再叠上分布校正的方差与截断偏差。奖励函数写得对不对、探索够不够、状态是不是部分可观测、上线之后安不安全，这些问题一个也没有因为多训了一个网络而消失。判断训练是否健康，至少要同时看策略熵有没有过早塌缩、回报分布里有没有极端值在主导梯度、Critic 的预测与实际回报差多少、优势的量级是否合理、重要性比率有没有长尾，以及不同种子之间的差距是否还在可接受范围。只盯平均回报那一根曲线，看不出 Critic 是不是在乱说，也排除不了策略正在钻奖励函数的漏洞。

回头看这一路：动态规划要求模型已知，换来对整个状态空间的精确保证；采样方法放弃模型，换来只在被访问过的地方有效；函数逼近放弃逐状态的独立表示，换来大空间可处理，同时交出了压缩性；离策略放弃数据与目标一致，换来数据可复用，交出了分布的对齐。每放松一次，适用范围扩大一圈，能给出的保证削掉一层。这个单元想留下的判断只有一条：知道自己此刻站在哪一层，手里还剩下哪些保证。
